% LaTeX template for AAGI Reports
% Date last modified: 07/11/2025

\documentclass[12pt]{article}

%% Configuration %%%%%%%%%%%%%%%%%%%%%%%%%%%%%%%%%%%%%%%%%%%%%%%%%%%%%%%%%%%%%%%
% (Modify these if you have a need to do so, or if you have additional
% packages that you need to load. Otherwise this configuration should 
% be left as-is.)
\usepackage{lmodern}
\usepackage{microtype}
\usepackage[paper=a4paper,top=35mm,left=25mm,right=25mm,bottom=30mm]{geometry}
\usepackage{adjustbox}
\usepackage[title]{appendix}
\usepackage{parskip}
\usepackage{amsmath}
\usepackage{amssymb}
\usepackage{titlesec}
\usepackage{fancyhdr}
\usepackage{fancyvrb}
\usepackage{rotating}
\renewcommand{\familydefault}{\sfdefault}

% AAGI colours
\usepackage[dvipsnames]{xcolor}
\usepackage{graphicx}
\definecolor{aagiteal}{HTML}{00808b}
\definecolor{aagigreen}{HTML}{b6d438}
\definecolor{aagiblack}{HTML}{414042}

% Command definitions and settings copied over from the old AAGI.sty
\newcommand{\HRule}{\hspace{-25mm}\rule{160mm}{0.6mm}}
\setlength{\parindent}{0cm}
\setlength{\headheight}{2pt}
\setlength{\parskip}{2.5ex}
\setlength{\emergencystretch}{3em}
\DeclareOldFontCommand{\sbf}{\normalfont\bfseries\sffamily}{\mathsf}
\DeclareOldFontCommand{\tbf}{\normalfont\bfseries\ttfamily}{\mathtt}
\pagestyle{fancy}
\renewcommand{\sectionmark}[1]{\markboth{\thesection\hspace{2.5mm} #1}{}}
\fancyhf{}
\fancyhead[L]{\sffamily\leftmark}
\fancyfoot[C]{\sffamily \thepage}
\fancypagestyle{bibliography}{\fancyhead[L]{\sffamily References}}

% Colour tables + long tables
\usepackage{colortbl}

\usepackage{longtable}
\setlength{\LTcapwidth}{\textwidth}

% Figures configuration
\usepackage{caption}
\captionsetup{font={sf}}
\setlength{\fboxsep}{0pt}

% URL styling
\usepackage{url}

% FontAwesome icons
\usepackage{fontawesome5}

% Tikzpictures/PGFplots (if desired)
%\usepackage{tikz}
%\usepackage{pgfplots}
%\pgfplotsset{compat=newest}

% Listings for code (if desired)
%\usepackage{inconsolata}
%\usepackage{listings}
%\lstset{%
%  basicstyle=\footnotesize\ttfamily,
%  keywordstyle=\footnotesize\ttfamily,
%  language=R}

% Bibliography (bibtex, natbib)
\usepackage{natbib}


% A fancy coloured box for notes/warnings
\usepackage[most]{tcolorbox}
\newtcolorbox{note}[1][]{%
  colback=aagiteal!7, colframe=aagiteal!7, notitle, sharp corners,
  borderline west={2pt}{0pt}{aagiteal}, borderline east={2pt}{0pt}{aagiteal},
  enhanced, breakable}

% Hyperref for links/cross-referencing
\usepackage{hyperref}
\hypersetup{%
  colorlinks=true, linkcolor=aagiteal, urlcolor=aagiteal, citecolor=aagiteal}
\def\sectionautorefname{Section}
\def\subsectionautorefname{Subsection}
\def\appendixautorefname{Appendix}
\def\figureautorefname{Figure}
\def\tableautorefname{Table}


%% Document %%%%%%%%%%%%%%%%%%%%%%%%%%%%%%%%%%%%%%%%%%%%%%%%%%%%%%%%%%%%%%%%%%%%
%EDIT these parameters for your project here
\newcommand{\myuniversity}{University Name}
\newcommand{\myaaginode}{AAGI-XX}
\newcommand{\myprogram}{AAGI Program}  % e.g. Service \& Support,
                                       %      Research \& Development,
                                       %      Upskilling \& Awareness

\begin{document}
  %% Title page
  \thispagestyle{empty}
  \vspace*{-30mm}
  \begin{center}
    \includegraphics[height=3cm]{assets/AAGI_logo.pdf}    
  \end{center}
  \vspace*{2cm}
  \begin{flushright}
    \HRule\\[5mm]
    \adjustbox{tabular=@{}r@{}r@{},width=\linewidth}{%
      \color{aagiteal}\sbf Analytics for the Australian Grains Industry\\[1pt]
      \color{aagiteal}\sbf {\myuniversity} ({\myaaginode})\\[0pt]
      \color{aagiteal}\sbf {\myprogram} Report\\}\\[0pt]
    \vspace*{6mm}
    {\LARGE\sbf%
       Informative title for report  %EDIT your report title here
    }\\[4pt]
    {\Large\sbf%
      Report for GRDC/AAGI Project Code  %EDIT your project code here
    }\\[8mm]
    
    %EDIT your report audience/writers here
    \sf\normalsize
    Prepared for:
      {\sbf Client Name (Client Organisation)}\\
      Name 2 (Organisation 2)  % and so on...
    
    Prepared by:
      {\sbf Your Name}\\
      Name 2\\
      Name 3  % and so on...

    %EDIT your email here
    \newcommand{\myemail}{your.email@example.edu.au}
    Correspondence email:\; \texttt{\href{mailto:{\myemail}}{\myemail}}
      
    \today\\[-2pt]
    \HRule
  \end{flushright}

  \vfill

  \noindent\makebox[\textwidth]{%
    \includegraphics[width=\paperwidth]{assets/Partners.pdf}}
  \vspace*{-36mm}
  
  \clearpage
  \pagestyle{plain}
  \setcounter{page}{1}
  \tableofcontents

  % Other lists of tables or figures (uncomment as desired)
  %\listoftables
  %\addcontentsline{toc}{chapter}{List of Tables}
  %\listoffigures
  %\addcontentsline{toc}{chapter}{List of Figures}

  %% Document proper
  \clearpage
  \phantomsection
  \section*{Executive Summary}
  \addcontentsline{toc}{section}{Executive summary}
  Executive summary goes here. Summarise the work that was provided by AAGI and
  the main results. (Try to limit this to a single page.)

  \pagestyle{fancy}

  \clearpage
\section{Introduction}
  
  The report body text goes here.
  
  Note that the following section titles are provided as a guide, but feel free
  to modify them in any way that you like if they are not relevant to your work
  of if using a different set of titles enhances the readability of your report.
  
  Use AAGI theme colours for emphasis (\verb|\color{aagiteal}|
  for {\color{aagiteal}AAGI teal}, \verb|\color{aagigreen}| for
  {\color{aagigreen}AAGI green}, and \verb|\color{aagiblack}| for
  {\color{aagiblack}AAGI black}).
  
  Use \verb|\citep{blah}| for parenthetical
  citations: \citep{blah}. Use \verb|\citet{blah}| for text citations, e.g.
  \citet{blah}. References should be entered in bibtex format in the
  \texttt{references.bib} file.
  
  A simple AAGI-themed table could be done as in \autoref{tab:aagi}.
  \begin{table}[h]
    \caption{An AAGI-themed table.}\label{tab:aagi}
    \centering
    \begin{tabular}{| l | r r |}
    \rowcolor{aagiteal} {\color{white}Table Heading} & %
    {\color{white}Next Heading} & {\color{white}Another Heading} \\
    \hline
    blah & 2 & 3 \\
    \hline
    \end{tabular}
  \end{table}
  
  Use \verb|\begin{note} ... \end{note}| for emphatic notes throughout
  the report if you like. You can use \verb|\faIcon{...}| to add 
  FontAwesome icons
  for decoration, as shown below.
  
  \begin{note}
    {\sbf\color{aagiteal!75!black} \faIcon{exclamation-triangle} NOTE}\quad
    This is an example of a note.
  \end{note}
  
  \vspace*{1cm}  
  
  In this introduction section, outline the goals of the research project and
  any necessary background to provide contexts and rationale for the
  research/work.

  \subsection{This is a subsection}
  
  Use subsections (and subsubsections) as much as you like to help structure
  your content.

  \subsubsection{This is a subsubsection}

  \clearpage
  \section{Trial design}
  
  Include the trial layout, type, and the experimental design, such as treatments,
  number of replicates, blocking, etc. Include graphics and pictures where useful.
  
  Make sure to include notes about any specific considerations that must be kept
  in mind for the type of trial (e.g., small plots, glasshouse experiments,
  genetics, breeding trials, OFE projects, bioinformatics projects, and so on).
  
  \clearpage
  \section{Exploratory data analysis}
  
  You can use a dedicated section here for the preliminary data exploration and
  dataset visualisation. This can be useful to aid in the interpretation of the
  data and graphs generated throughout the work, as well as for providing a
  rationale for the statistical methodologies used.  
  
  \clearpage
  \section{Methods}
  
  Use this section to include a detailed description of the procedures and
  methodologies used for the analysis. Include versions/commits on developed
  pipelines, code scripts, and input/output details as relevant.
  
  Detail also the approach taken for the data analysis work here (or in a
  separate section).
  
  \clearpage
  \section{Results and discussion}
  
  Include all of the results of the analytical work, and discuss the findings
  and their implications with respect to the research question.
  
  \clearpage
  \section{Datasets and metadata}
  
  (Optional) Where useful, include descriptions of the datasets relevant to
  this research, including links and DOIs to enable findability of these datasets.
  (Often for GRDC projects these will have associated metadata entries in the GRDC
  Data Catalogue, which can be linked where appropriate.)
  
  Where you have publicly available code on GitHub associated with the analytics
  (e.g., pipelines), you can include git commit numbers and tags. Including
  md5/sha256 checksums for input data and outputs is also useful for
  research reproducibility.
  
  \clearpage
  \section{Trial location maps}
  
  (Optional) Include if relevant to the project.
    
  
  % Appendices
  \clearpage
  \begin{appendices}
    \clearpage
    \section{This is an appendix}
    
    Include Appendices at the end of the report as necessary, to contain any
    additional supporting information of interest.
    
    \subsection{This is an appendix subsection}

    \clearpage
    \section{This is another appendix}
    And so on. Include as many appendices as you like.    
    
  \end{appendices}

  % Bibliography
  \clearpage
  \phantomsection
  \pagestyle{bibliography}
  \addcontentsline{toc}{section}{References}
  \bibliographystyle{assets/biom}
  \bibliography{references.bib}
\end{document}
